\PassOptionsToPackage{hyphens}{url}
\documentclass[10pt]{article}

\usepackage[english]{babel}
\usepackage[T1]{fontenc}
\usepackage[utf8]{inputenc}
\usepackage{textcomp}
\usepackage[letterpaper,top=2.2cm,bottom=1.7cm,left=1.6cm,right=1.6cm,marginparwidth=1.75cm]{geometry}
\usepackage{amssymb}
\usepackage{booktabs}
\usepackage{array}
\usepackage{graphicx}
\usepackage{xcolor}
\usepackage[colorlinks=true, allcolors=blue]{hyperref}
\usepackage{enumitem}
\usepackage{float}
\usepackage{fancyhdr}
\usepackage{lastpage}
\usepackage{titlesec}

\setlist{nosep,leftmargin=1.1em}
\titleformat{\section}{\normalsize\bfseries}{\thesection.}{0.4em}{}
\titlespacing*{\section}{0pt}{4pt}{1pt}
\setlength{\parindent}{0pt}
\setlength{\parskip}{2pt}

% Red text = fill in / verify before submitting
\newcommand{\todo}[1]{\textcolor{red}{[#1]}}

\pagestyle{fancy}
\fancyhf{}

\setlength{\headheight}{0.8cm}
\setlength{\footskip}{1.1cm}

\newcommand{\groupnames}{\begin{tabular}[t]{@{}l@{}}
GUERRAS Samy Cherif
\end{tabular}}

\fancyhead[L]{\includegraphics[width=1.5cm]{University_of_Luxembourg_logo_(fr).svg.png}}
\fancyhead[C]{\textbf{M2 -- Technology Exploration and Proposal}}
\fancyhead[R]{\footnotesize Digital Transformation (BPINFOR-131)}

\fancyfoot[L]{\footnotesize\groupnames}
\fancyfoot[R]{Page \thepage\ of \pageref{LastPage}}

\fancypagestyle{plain}{
  \fancyhf{}
  \fancyhead[L]{\includegraphics[width=1.5cm]{University_of_Luxembourg_logo_(fr).svg.png}}
  \fancyhead[C]{\textbf{M2 -- Technology Exploration and Proposal}}
  \fancyhead[R]{\footnotesize Digital Transformation (BPINFOR-131)}
  \fancyfoot[L]{\footnotesize\groupnames}
  \fancyfoot[R]{Page \thepage\ of \pageref{LastPage}}
}

\begin{document}
\small\sloppy
\textbf{Team:} Markr \quad \textbf{Lead:} Guerras Samy \quad \textbf{Product:} Markr \quad \textbf{Repo:} \url{https://github.com/Guerrassc/Markr} (tag \texttt{m2}, public)

\section{Approaches compared}
Task: turn an assignment brief into a complete checklist of requirements. Same 10 briefs for every approach (\texttt{experiments/m2\_extraction}). Cost uses the price list in section 5.
\footnotesize
\begin{tabular}{@{}p{2.9cm}p{3.9cm}p{4.5cm}p{4.5cm}@{}}
\toprule
 & \textbf{A: regex rules (baseline)} & \textbf{B: small LLM} (gemini-3.5-flash-lite) & \textbf{C: large LLM} (gemini-3.8-flash) \\
\midrule
Output quality (recall) & 50\,\% (measured) & \todo{from results} & \todo{from results} \\
Cost per user action & \texttt{\texteuro{}}0 & $\approx$\$0.0033 & $\approx$\$0.0060 \\
Latency (median / worst) & under 1 ms & \todo{s / s} & \todo{s / s} \\
Data needed & none & none (prompt only) & none (prompt only) \\
Failure modes & misses rules without a date or number (Harvard, plain text, ``every member must speak'') & \todo{from failed cases} & \todo{from failed cases} \\
Lock-in / switching & none & low: prompt in \texttt{prompts/}, plain HTTP client, OpenAI-compatible API & low (same) \\
\bottomrule
\end{tabular}
\small

\section{Experiment}
\texttt{experiments/m2\_extraction} runs all three approaches on 10 synthetic briefs (essay, lab, code, rubric, French, informal email, \dots) with labelled requirements. Metric: recall (a requirement counts as found when one extracted line contains all its key words). How to run: that folder's README; raw output: \texttt{results/m2\_results.md}.\\
\textbf{Result:} regex baseline 50\,\% mean recall at 3.7 items per brief (best: thesis 71\,\%, worst: essay 29\,\%). LLM rows: \todo{paste table}. \textbf{Commentary:} \todo{what the LLMs found that rules missed; two or three concrete missed items; any failures}.

\section{Chosen approach, limitations, switch trigger}
\textbf{Chosen:} \todo{gemini-3.8-flash or gemini-3.5-flash-lite} through Google\'s Gemini API, prompt kept in the repo, \todo{reason from the results}. \textbf{Limitations:} (1) only 10 synthetic briefs and keyword matching, so recall is optimistic and can miss correct paraphrases; (2) briefs and drafts are student work sent to Google, and on the free tier Google may use content to improve its products, so testing uses only synthetic briefs and real use needs the paid tier; (3) scanned or oddly laid-out PDFs need text extraction first and may lose rules. \textbf{We would switch} to the small model if its recall is within 5 points of the large one, and to another provider if rate limits or terms block the MVP demo.

\section{First architecture sketch}
Student $\rightarrow$ web UI (brief + draft) $\rightarrow$ Markr backend $\rightarrow$ \textbf{LLM API} (Google Gemini API, \todo{model}) $\rightarrow$ checklist and check result $\rightarrow$ UI. Prompt: \texttt{prompts/extract\_requirements.txt}. The API key lives in a git-ignored \texttt{.env}. Brief and draft text leave our system for the API provider. Mermaid sketch: \texttt{docs/architecture.md}.

\section{First cost estimate}
Google Gemini API price list, paid tier, checked 19 Oct 2026 (ai.google.dev/gemini-api/docs/pricing), per 1M tokens: gemini-3.8-flash \$0.75 input and \$3.75 output (rising to \$1.50 and \$7.50 on 1 Jan 2027); gemini-3.5-flash-lite \$0.30 and \$2.50. Output includes thinking tokens. One action = extract the checklist from one brief and check one draft. Assumption: 3,800 input and 850 output tokens (brief about 700 tokens, draft about 2,800, prompts and checklist the rest); \todo{replace with measured tokens from \texttt{results/m2\_results.json}}.\\
Large: (3,800 $\times$ 0.75 + 850 $\times$ 3.75) / 1,000,000 = 0.00285 + 0.00319 = \textbf{\$0.00604} per action. Small: (3,800 $\times$ 0.30 + 850 $\times$ 2.50) / 1,000,000 = 0.00114 + 0.00213 = \$0.00327. At 3,000 actions per month: about \$18.11 (large) or \$9.80 (small); after the 1 Jan 2027 price rise the large model costs about \$0.0121 per action.

\section{Plan to MVP (due Mon 16 Nov)}
\footnotesize
\begin{tabular}{@{}p{1.5cm}p{8.2cm}p{1.2cm}p{4.2cm}@{}}
\toprule
\textbf{Date} & \textbf{Task} & \textbf{Owner} & \textbf{Done when} \\
\midrule
20 Oct & Run all 3 approaches, pick the model, commit results & Samy & table in \texttt{results/} \\
24--28 Oct & PDF/text upload with text extraction; checklist extraction end to end with the chosen model & Samy & 5 sample briefs, no mocked responses \\
1--5 Nov & Draft-versus-checklist check; simple web UI joining both steps & Samy & full flow in a browser \\
9 Nov & Test on 5 real briefs (with consent); log bugs & Samy & 3+ bugs listed \\
12--14 Nov & README setup in 5 steps, architecture diagram, run on another laptop, 3-minute demo video & Samy & clean machine works; unlisted video link \\
\bottomrule
\end{tabular}
\small
Repository created, README skeleton in place, public (no need to add the lecturers).

\end{document}
